\documentclass[10pt,a4paper]{amsart}
\usepackage[margin=19mm]{geometry}
\usepackage{amsmath,amssymb,mathtools}
\usepackage[scr=rsfs,cal=euler]{mathalfa}
\usepackage{xfrac}
\usepackage[unicode,hidelinks,bookmarks=false]{hyperref}
\newcommand{\ie}{i.e.\ }
\newcommand{\cf}{cf.\ }
\newcommand{\resp}{respectively\ }
\newcommand{\Subsection}{\subsection}
\providecommand{\nobreakdash}{\nobreak\hbox{-}}
\setlength{\emergencystretch}{2em}
\multlinegap0pt
\usepackage[shortlabels]{enumitem}
\usepackage{mathtools}
\usepackage{amssymb}
\usepackage{esint}
\usepackage{dsfont}
\usepackage{xcolor}
\newtheorem{theorem}{Theorem}
\newcommand{\Hil}{\mathscr{H}}
\newcommand{\R}{\mathbb{R}}
\newcommand{\dom}{\mathrm{dom}}
\newcommand{\embed}{\hookrightarrow}% Embeded in
\renewcommand{\epsilon}{\varepsilon} 
\title{First-order elliptic boundary value problems on manifolds with noncompact boundary}
\author{Christian Bär, Lashi Bandara}
\date{Source: arXiv:2401.17784v3 (CC BY 4.0)}
\begin{document}
\maketitle

\noindent\emph{Source and licence.} First-order elliptic boundary value problems on manifolds with noncompact boundary. Version arXiv:2401.17784v3 (CC BY 4.0), distributed by arXiv under the Creative Commons Attribution 4.0 International licence, http://creativecommons.org/licenses/by/4.0/, as recorded in the arXiv licence metadata for this exact version and shown on the abstract page https://arxiv.org/abs/2401.17784v3. Full licence notice: \texttt{licenses/arxiv-2401.17784-CC-BY-4.0.txt}.\par\smallskip

\noindent\textit{Editorial scope.} Counted unit: the article's theorem Thm:AbsCpctEmbed (source0.tex lines 3382--3394). The article's complete original proof of it is delivered with it (source0.tex lines 3396--3474, one \texttt{proof} environment of 79 lines). No mathematics is written, repaired or re-derived: the statement and the proof are the article's own lines, in the article's own notation, and no retained line is substituted or re-flowed.  No further source-local context is needed: every cross-reference of every retained line resolves inside the two delivered spans.  Nothing was omitted from any retained span, so the package carries no omission record.  No retained line contains a citation, so the wrapper carries no bibliography and no bibliography entry.  No retained line contains a figure environment, an image-inclusion command or drawing code, so no image is delivered and the package declares no asset dependency.  Editorial formatting only, in addition: an AMS \texttt{amsart} preamble that restates the article's own declarations and macros used by the retained lines, namely the environments \texttt{theorem} on separate counters, together with the standard packages those lines need.  The retained environments print in this document as Theorem 1 in order of appearance, whereas the article prints the same items with the numbering of its own full document.  The complete native source of the article as distributed in the arXiv e-print for this exact version is bundled unchanged as \texttt{sources/153804/source0.tex}.
\begin{theorem}
\label{Thm:AbsCpctEmbed}
Let $T$ be a selfadjoint operator on a separable Hilbert space $\Hil$.
Then the following are equivalent.
\begin{enumerate}[label=(\roman*), labelwidth=0pt, labelindent=2pt, leftmargin=21pt]
\item \label{It:AbsCpctEmbed:1}
$T$ has discrete spectrum.
\item \label{It:AbsCpctEmbed:2} 
The embedding $\dom((1+T^2)^s) \embed \dom((1+T^2)^t)$ is compact for some $s, t \in \R$ with $s > t\ge0$.
\item  \label{It:AbsCpctEmbed:3} 
The embedding $\dom((1+T^2)^s) \embed \dom((1+T^2)^t)$ is compact for all $s, t \in \R$ with $s > t\ge0$.
\end{enumerate}
\end{theorem}

\begin{proof}
Without loss of generality, we assume that $T$ unbounded.
Indeed, if $T$ is bounded, then each of the assertions \ref{It:AbsCpctEmbed:1}--\ref{It:AbsCpctEmbed:3} is equivalent to $\Hil$ being finite dimensional.

It is immediate that \ref{It:AbsCpctEmbed:3} implies \ref{It:AbsCpctEmbed:2}.
We show that \ref{It:AbsCpctEmbed:2} implies \ref{It:AbsCpctEmbed:1}.
So, assume \ref{It:AbsCpctEmbed:2} and note that 
$$ 
\dom((1+T^2)^s) \xhookrightarrow{\mathrm{compact}} \dom((1+T^2)^t) \xhookrightarrow{\mathrm{continuous}} \Hil,
$$
and therefore we have that $(1+T^2)^{-s}:\Hil \to \Hil$ is a compact map.
Hence $(1+T^2)^{-s}$ has discrete spectrum except for the accumulation point $0$.
Since $s>0$ it follows that $1+T^2$ and hence $T$ has discrete spectrum.

It remains to show that \ref{It:AbsCpctEmbed:1} implies \ref{It:AbsCpctEmbed:3}.
Let $T$ have discrete spectrum.
To show that $\dom((1+T^2)^s) \embed \dom((1+T^2)^t)$ is compact, let $u_k\in \dom((1+T^2)^s)$ be a bounded sequence, 
$$
\|u_k\|_{\dom((1 +T^2)^s)} \leq C.
$$
We need to find a subsequence which converges in $\dom((1 +T^2)^t)$.
Let $\lambda_j$ be the eigenvalues of $T$ and $\phi_j$ its orthonormalised eigenvectors.
We write $u_k = \sum_j u_{k,j}\phi_j$.
For each fixed $j$ we have
$$
(1+\lambda_j^2)^{2s} |u_{k,j}|^2 \le \|u_k\|_{\dom((1 +T^2)^s)}^2 \le C^2
$$
and hence 
$$
|u_{k,j}| \le C .
$$
By a diagonal argument, we can pass to a subsequence, again denoted $(u_k)$, such that for each $j$
$$
u_{k,j} \xrightarrow{k\to\infty} w_j .
$$

We show that $(u_k)$ is a Cauchy sequence in $\dom((1 +T^2)^t)$.
Let $\epsilon>0$.
Since the spectrum is discrete and $T$ is unbounded there exists $j_0$ such that 
$$
|\lambda_j| \ge \frac1\epsilon 
$$
for all $j> j_0$.
For these $j$ we then have
$$
(1+\lambda_j^2)^{t-s}
\le 
|\lambda_j|^{2(t-s)}
\le
\epsilon^{2(s-t)}.
$$
Therefore,
\begin{align*}
\sum_{j> j_0} (1+\lambda_j^2)^{2t} |u_{k,j}-u_{\ell,j}|^2
&=
\sum_{j> j_0} (1+\lambda_j^2)^{2(t-s)}(1+\lambda_j^2)^{2s} |u_{k,j}-u_{\ell,j}|^2 \\
&\le
\epsilon^{4(s-t)}\sum_{j\le j_0}(1+\lambda_j^2)^{2s} |u_{k,j}-u_{\ell,j}|^2 \\
&\le
\epsilon^{4(s-t)} \|u_k-u_\ell\|_{\dom((1 +T^2)^s)}^2 \\
&\le
\epsilon^{4(s-t)} \cdot 4C^2.
\end{align*}
Choose $m$ so large that such that we have for all (finitely many) $j\le j_0$
$$
(1+\lambda_j^2)^{2t} |u_{k,j}-u_{\ell,j}|^2 \le \frac{\epsilon}{j_0}
$$
whenever $k,\ell\ge m$.
Then we find for such $k$ and $\ell$ that
\begin{align*}
\|u_k-u_\ell\|_{\dom((1 +T^2)^t)}^2
&=
\sum_{j\le j_0}(1+\lambda_j^2)^{2t} |u_{k,j}-u_{\ell,j}|^2
+ \sum_{j> j_0}(1+\lambda_j^2)^{2t} |u_{k,j}-u_{\ell,j}|^2 \\
&\le
\epsilon + \epsilon^{4(s-t)} \cdot 4C^2.
\end{align*}
This shows that $(u_k)$ is a Cauchy sequence in $\dom((1 +T^2)^{\frac t 2})$ and hence converges.
\end{proof} 

\end{document}
